\section{Attention as a Proxy}
\label{sec:proxy}

Attention-based hallucination methods share a convenient but under-tested proxy: they treat where an LVLM routes computation as evidence for whether an object claim is visually supported. Detection methods instantiate this proxy by converting attention into a scalar score. SVAR sums object-to-image attention mass \cite{jiang2024devils}; PAS measures object-token attention to previously generated prelim tokens \cite{hoangxuan2026pas}; and recent fine-grained detectors combine patch attention dispersion with cross-modal consistency \cite{beyond2026global}. These methods are attractive because they are training-free and internally available, but their scores still summarize routing rather than verification.

Mitigation methods instantiate the same proxy operationally. PAI increases image-token attention logits and combines this with contrastive/CFG decoding in the full method \cite{liu2024pai}; ClearSight scales mid-layer text-to-image attention and suppresses system-prefix attention \cite{yin2025clearsight}; Visual Attention Sink redistributes mass from high-activation sink tokens toward non-sink visual tokens \cite{kang2025visualsink}. All three are motivated by the idea that increasing useful visual attention should improve grounding. The recent ``Mirage'' analysis of contrastive decoding shows that POPE gains can instead come from output-distribution shifts and decoding constraints rather than hallucination mitigation \cite{yin2025mirage}. We ask whether attention-side methods suffer an analogous failure: their routing changes may be real, while their grounding gains are not.

\begin{table}[t]
\centering
\scriptsize
\setlength{\tabcolsep}{3.0pt}
\begin{tabular}{p{0.23\columnwidth}p{0.32\columnwidth}p{0.32\columnwidth}}
\toprule
Use & Proxy & Failure mode tested \\
\midrule
Detection & attention mass or shape is a grounding score & score tracks generation position or object frequency \\
Mitigation & more visual routing reduces hallucination & output prior or caption style shifts without verification \\
Stress test & purified shape removes mass bias & normalized shape still tracks decoding dynamics \\
\bottomrule
\end{tabular}
\caption{The shared proxy behind the methods we test. The paper evaluates whether the proxy survives controls that target its most likely failure mode.}
\label{tab:proxy}
\end{table}



The missing distinction is between \emph{allocation} and \emph{verification}. Allocation says that a token attends to the image, to prior text, or to a sink position. Verification asks whether the image contains the object being claimed. A method that only changes allocation can increase visual plausibility, object richness, or the probability of answering ``yes'' without checking object presence. This distinction motivates our diagnostic design: every score or intervention is evaluated not only by its aggregate metric, but also by whether it remains selective after controlling for nuisance variables and output-prior shifts.

\paragraph{Associated evidence is not target evidence.} The allocation--verification gap is sharpest when the attended region is meaningful but non-discriminative. In many hallucinations, the visual evidence is not arbitrary noise: it is a nearby object, a scene context, or a visual token whose language-space neighborhood is associated with the target. A bathroom region may support the prior that a sink is plausible, but it is not evidence that a sink is present. This is precisely the case where attention explanations are most tempting and most dangerous: the heatmap looks visually sensible, while the object-level claim is false. We therefore treat \emph{target-discriminative selectivity} as the missing criterion. A useful detector should separate present and absent instances of the same object under comparable position and context; a useful intervention should increase grounding for true positives more than for false positives.

This framing also explains why a negative result on simple attention scores is still informative. Attention may participate in the computation that supports grounding, but a scalar proxy must be stable to nuisance variables and selective for object evidence. If a score only becomes strong when late objects are more likely to hallucinate, or if an intervention increases visual routing for both true and false positives, then the method has not measured verification. Our experiments therefore focus less on whether attention changes and more on whether the change is \emph{selective} in the direction that grounding requires.
